\documentclass[UTF8,12pt,a4paper]{ctexart}

\usepackage{amsmath}
\usepackage{cases}
\usepackage{cite}
\usepackage{graphicx}
\usepackage{enumerate}
\usepackage{algorithm}
\usepackage{caption}   % \caption*需要
\usepackage{subcaption} % 子图布局
\usepackage[noend]{algpseudocode} %algorithmicx
\makeatletter
\renewcommand{\fnum@algorithm}{\fname@algorithm}
\makeatother
\renewcommand{\algorithmicrequire}{\textbf{Input:}}
\renewcommand{\algorithmicensure}{\textbf{Output:}}
\usepackage[margin=1in]{geometry}
\geometry{a4paper}
\usepackage{fancyhdr}
\pagestyle{fancy}
\fancyhf{}
\usepackage{xcolor}

% ------------------------- 代码展示设置 -------------------------
\usepackage{listings}
\lstset{
	basicstyle=\ttfamily,                                % 传统上，代码打字机字体好些
	breaklines,
	tabsize=4,
	extendedchars=false,
	columns=fixed,
	%numbers=left,                                        % 在左侧显示行号
	numbers=none,
	frame=none,                                          % 不显示背景边框
	backgroundcolor=\color[RGB]{245,245,244},            % 设定背景颜色
	keywordstyle=\color[RGB]{40,40,255},                 % 设定关键字颜色
	numberstyle=\footnotesize\color{darkgray},           % 设定行号格式
	commentstyle=\it\color[RGB]{0,96,96},                % 设置代码注释的格式
	stringstyle=\rmfamily\slshape\color[RGB]{128,0,0},   % 设置字符串格式
	showstringspaces=false,                              % 不显示字符串中的空格
	xleftmargin=1em,
	%xrightmargin=1em,
	%aboveskip=-0.5em
	language=c++,                                        % 设置语言
}



% ------------------- 设置超链接，便于跳转 -----------------
\usepackage{enumitem}
\usepackage{hyperref}
\hypersetup{
    pdfborder={0 0 0},    % 消除超链接边框
    colorlinks=true,       % 将边框改为颜色标记（可选）
    linkcolor=black,       % 内部链接颜色（目录、引用等）
    citecolor=black,       % 引用颜色
    urlcolor=blue          % URL链接颜色（保持蓝色以便区分）
}




% ------------------------ 标题区 ----------------------------
\title{算法原理 HW1/Lab1 作业/实验报告模板}
\author{
	姓名：\underline{你的名字}~~~~~~
	学号：\underline{你的学号}~~~~~~}
\date{\today}
\pagenumbering{arabic}

\begin{document}



% ------------------------ 页眉和页脚 ----------------------------
\fancyhead[L]{你的名字}
\fancyhead[C]{分治算法}
\fancyfoot[C]{\thepage}

\maketitle
%\tableofcontents
%\newpage




% ------------------------- 伪代码展示 ----------------------------
\begin{algorithm}[H]%
	\caption{RAND-MIN-CUT($G=(V,E)$)}
	\begin{algorithmic}[1]
		\State 为图$G$的所有边都随机赋予一个正权值
		\State 找出$G$的最小生成树$T$
		\State 删除$T$中权值最大的一条边得到两棵子树$T_1,~T_2$
		\State 令$T_1$的顶点集为$V_1$，则$T_2$的顶点集为$V_2=V-V_1$
		\State $C=\{ uv|uv\in E, u\in V_1, v\in V_2 \}$
		\State \Return{$C$}
	\end{algorithmic}
\end{algorithm}

上述算法中，第1行XXX；第XXX行开始循环，处理XXX。算法的时间复杂度为$O(n^2)$。



% ------------------------- 代码展示 ----------------------------
\begin{lstlisting}[language={C++}, basicstyle=\ttfamily\small]
	main() {
		int n, total, x, y, z;
		scanf("%d",&n);
		total = 3;
		while(1){
			for(x=1; x<=total-2; x++)
			for(y=1; y<=total-x-1; y++){
				z = total - x - y;
				if( exp(x,n) + exp(y,n) == exp(z,n))
				printf("hello, world\n");
			}
			total++;
		}
	}
\end{lstlisting}


% ------------------------- 表格 ----------------------------
\begin{table}[h]
    \centering
    \begin{tabular}{|c|c|c|c|c|c|}
    \hline
    数据集规模 & 100 & 1,000 & 10,000 & 100,000 & 1,000,000 \\
    \hline
    暴力搜索改良版 & 0.0443 & 0.9452 & 42.6917 & 4115 & 401570\\
    \hline
    先排序后扫描 & 0.0152 & 0.2597 & 2.4512 & 29.8881 & 381.301\\
    \hline
    分治算法 & 0.0731 & 0.6339 & 5.2765 & 52.9706 & 541.206\\
    \hline
    哈希表方法 & 0.003 & 0.0078 & 0.0422 & 0.3377 & 3.1189\\
    \hline
\end{tabular}
\caption{五种算法在不同规模数据集下的运行时间测试(时间单位: ms)}
\label{tab:time_results}
\end{table}


% ------------------------- 复杂表格 ----------------------------
\setlength{\arrayrulewidth}{0.9pt} % 线条粗一些，可按需调
\begin{table}[h]
    \centering
    % 如果过宽，可取消注释用 \resizebox 缩放
    \resizebox{\textwidth}{!}{%
    \begin{tabular}{|P{2.6cm}|P{2.6cm}|P{4.0cm}|P{2.2cm}|P{2.2cm}|P{2.4cm}|P{1.8cm}|}
        \hline
        % ----------------------- 表头 ------------------------
        \makecell{Process} & Channel & \makecell{Width/Width$_{\text{eff}}$(nm)} &
        $I_{\text{on}}$ & $I_{\text{off}}$ & $I_{\text{on}}/I_{\text{off}}$ & $V_{\text{T}}$ \\
        % ----------------------- 主体内容 ------------------------
        \cline{1-7}
        \multirow{2}{*}{Bulk}   & NMOS & 128 & 168.3$\mu$A & 185nA & 910 & 0.1377V \\
        \cline{2-7}
                                & PMOS & 128 & 92.59$\mu$A & 204.7nA & 452 & -0.1333V \\
        \cline{1-7}
        \multirow{2}{*}{FinFET} & NMOS & \makecell{128 (Fin num=2)} & 157.3$\mu$A & 89.76nA & 1752 & 0.1246V \\
        \cline{2-7}
                                & PMOS & \makecell{128 (Fin num=2)} & 141.4$\mu$A & 127.5nA & 1109 & -0.1120V \\
        \hline
    \end{tabular}}
    \caption{$I_{\text{on}}$, $I_{\text{off}}$ and $V_{\text{T}}$ for four transistors at temp = $90^\circ$C}
    \label{I_V_temp90}
\end{table}

% ------------------------- 图片 ----------------------------
\begin{figure}[htbp]
    \center
    \includegraphics[width=0.8\textwidth]{figure/wv_FinFET_25_90_off.png}
    \caption{$I_{\text{off}}$ comparison for PFET at $25^\circ$C and $90^\circ$C}
    \label{fig:Ioff_comparison}
\end{figure}



% -------------------------- 双图 ---------------------------
\begin{figure}[htbp]
    \centering
    % 第一张子图
    \begin{subfigure}[b]{0.45\textwidth} % 子图宽度为总宽度的45%
        \centering
        \includegraphics[width=\textwidth]{figure/wv_Task1_Fin_NFET.png} 
        \caption{NMOS: $I_{\text{DS}}-V_{\text{DS}}$}
        \label{fig:NFET_I_V}
    \end{subfigure}
    \hfill % 添加水平间距
    % 第二张子图
    \begin{subfigure}[b]{0.45\textwidth}
        \centering
        \includegraphics[width=\textwidth]{figure/wv_Task1_Fin_PFET.png} 
        \caption{PMOS: $I_{\text{DS}}-V_{\text{DS}}$}
        \label{fig:PFET_I_V}
    \end{subfigure}
    % 总图标题
    \caption{$I_{\text{DS}}-V_{\text{DS}}$ curve}
    \label{fig:FET_IV}
\end{figure}

% 参考文献如不需要，将如下两行注释掉
\bibliographystyle{plain}
\bibliography{./template}  %bib文件名

\end{document}
